% ===================================================================
%  اللوحة رقم ٥ — البيت وكيف يُبنى.
%  Traduction 2026 d'apres texto/io, controlee sur texto/fr, texto/en
%  et texto/es. Consigne : voir texto/ar/10-lawha-01.tex.
%
%  Le tableau 5 est un DIALOGUE, et les deux volumes composent le nom
%  du parleur en petites capitales : on garde \textsc{}, que le rendu
%  laisse tomber en gardant le texte.
%
%  Les renvois du plan sont des LETTRES — (a) le couloir, (b) le
%  salon... — et non des numeros. On les ecrit comme l'ido les ecrit,
%  y compris la ou l'imprimeur a compose « (1) » pour « (l) ».
% ===================================================================

%%K t05-apar-1 apar
\VUsaut{6.0mm}
\VUtitre{15.0pt}{60}{اللوحة رقم ٥}
\VUsaut{1.4mm}
\VUtitre{11.4pt}{0}{البيت وكيف يُبنى.}
\VUsaut{2.2mm}

%%K t05-tit-1 sub
\VUpk{10.2pt}{40}{لقاء موفَّق.}
\VUsaut{3.0mm}

%%K t05-01-1 p
\textsc{يوحنا}. --- يا عمّي، بودّي كثيرًا أن أعرف كيف يُبنى \VUgras{البيت}.

%%K t05-01-2 p
\textsc{العم} \textsuperscript{(80)}. --- هذا هيّن يا بنيّ؛ تعال معي أُرِك
البيت الذي يبنيه السيد برنار \textsuperscript{(79)} والذي سأسكنه.

%%K t05-01-3 p
\textsc{يوحنا}. --- ومن رسم \VUgras{مخطط} \textsuperscript{(76)} هذا البيت؟

%%K t05-01-4 p
\textsc{العم}. --- \VUgras{المهندس المعماري} \textsuperscript{(75)}؛ وضع رسمًا
واضحًا جدًا، فأُقِرّ، وبدأ \VUgras{المقاول} \textsuperscript{(77)} العمل.

%%K t05-01-5 p
\textsc{يوحنا}. --- ومن المقاول؟

%%K t05-01-6 p
\textsc{العم}. --- السيد بيضون، والد صاحبك يعقوب. وهو يدير العمال مع
السيد رو المهندس.

%%K t05-01-7 p
وها هو ذا! السيد بيضون يأتي في وقته تمامًا و\VUgras{مسطرته}
\textsuperscript{(78)} معه، والسيد رو ومعه مخططه.

%%K t05-01-8 p
صباح الخير يا سيديّ. كيف حالكما؟

%%K t05-01-9 p
\textsc{السيد رو}. --- بخير، شكرًا. صباح الخير يا يوحنا؛ ما أسرع ما كبر
الولد!

%%K t05-01-10 p
\textsc{العم}. --- حقًا، صار رجلًا صغيرًا. وهو جاد جدًا؛ يريد أن يُشرَح له
كل ما يراه. واليوم يودّ أن يعرف كيف يُبنى البيت.

%%K t05-tit-2 sub
\VUpk{10.2pt}{40}{العمال.}
\VUsaut{3.0mm}

%%K t05-01-11 p
\textsc{السيد بيضون}. --- فتعال معي إذًا يا صديقي الصغير، وأشرح لك في
الحال، فإني أحب الأطفال الراغبين في التعلم.

%%K t05-01-12 p
ترى ذلك العامل الذي في يده \VUgras{مِجرفة} \textsuperscript{(82)}: إنه
\VUgras{حفّار} \textsuperscript{(81)}. يحفر \VUgras{خندقًا} \textsuperscript{(46)}
لمدّ أنبوب الماء. وهو الذي حفر الأرض التي يقوم عليها البيت ليتمكن البنّاء
من إقامة \VUgras{الجدران} الأربعة. والجزء المدفون من تلك الجدران يسمى
\VUgras{الأساس} \textsuperscript{(2)}. والفراغ الذي يحصره الأساس يؤلف
\VUgras{القبو} \textsuperscript{(3)}. ولهذا القبو نافذة صغيرة تسمى
\VUgras{المنفس} \textsuperscript{(5)}، و\VUgras{باب} \textsuperscript{(4)}، ودرج.

%%K t05-01-13 p
ومضى \VUgras{البنّاء} \textsuperscript{(108)} يرفع الجدران تاركًا فتحات
لـ\VUgras{الأبواب} \textsuperscript{(8)} و\VUgras{النوافذ} \textsuperscript{(17)}.

%%K t05-01-14 p
\textsc{يوحنا}. --- لكني لا أراه، هذا البنّاء!

%%K t05-01-15 p
\textsc{السيد بيضون}. --- لقد فرغ من العمل في البيت الآن. إنه هناك عند
\VUgras{الإسطبل} \textsuperscript{(54)}، فوق \VUgras{سقالته} \textsuperscript{(110)}،
يبني جدار \VUgras{غرفة الغسيل} \textsuperscript{(56)}.

%%K t05-01-16 p
معه مالج وحوض و\VUgras{شاقول} \textsuperscript{(109)}. لكنك لن تحفظ هذه
الأسماء.

%%K t05-01-17 p
\textsc{يوحنا}. --- بل سأحفظها، فسأطلب من عمي مالجًا صغيرًا وحوضًا، وأصنع
الشاقول بنفسي من قطعة خيط.

%%K t05-tit-3 sub
\VUsaut{3.0mm}
\VUpk{10.2pt}{40}{المواد.}
\VUsaut{3.0mm}

%%K t05-01-18 p
\textsc{العم}. --- حسن جدًا. لكن قل لي يا يوحنا، ماذا يلزم لبناء بيت؟

%%K t05-01-19 p
\textsc{يوحنا}. --- يلزم \VUgras{حجر منحوت} \textsuperscript{(59)}
و\VUgras{مِلاط} \textsuperscript{(65)}.

%%K t05-01-20 p
\textsc{العم}. --- تمامًا. انظر إلى ذلك الرجل ذي القبعة الكبيرة على
\VUgras{عربته} \textsuperscript{(136)}. إنه \VUgras{عربجي} \textsuperscript{(135)}.
يجلب إلى هنا كل يوم \VUgras{كتل حجر} \textsuperscript{(60)}. يفرغ عربته في
الفناء، فينحت \VUgras{النحّاتون} \textsuperscript{(83)} الكتل وينشرونها
بـ\VUgras{منشار الحجر} \textsuperscript{(85)}. ويحملها \VUgras{عامل}
\textsuperscript{(146)} إلى البنّاء، فيضعها في موضعها.

%%K t05-01-21 p
وفي هذه الأثناء يهيّئ \VUgras{عاجن المِلاط} \textsuperscript{(87)} المِلاط.
يلزمه \VUgras{رمل} \textsuperscript{(62)} و\VUgras{جير} \textsuperscript{(63)}
و\VUgras{ماء} \textsuperscript{(64)}. والجير إلى جانبه في \VUgras{كيس}
\textsuperscript{(71)}؛ و\VUgras{مساعد البنّاء} \textsuperscript{(111)} يجلب له الرمل
في \VUgras{عربة يد} \textsuperscript{(112)}، و\VUgras{الصبي} \textsuperscript{(113)}
عند المضخة \textsuperscript{(58)}. وهو يملأ \VUgras{دلوًا} \textsuperscript{(114)}.
وسيكون المِلاط جاهزًا بعد قليل. فيأخذه \VUgras{حامل المِلاط}
\textsuperscript{(107)} ويصعد به السلّم إلى السقالة حيث يُتم بنّاء ذلك الجانب
من البيت.

%%K t05-01-22 p
والآن، ماذا تسمي العامل الذي يصنع \VUgras{السطح} \textsuperscript{(31)}؟

%%K t05-01-23 p
\textsc{يوحنا}. --- إنه \VUgras{المُقرمِد} \textsuperscript{(118)}.

%%K t05-tit-4 sub
\VUsaut{3.0mm}
\VUpk{10.2pt}{40}{سطح البيت.}
\VUsaut{3.0mm}

%%K t05-01-24 p
\textsc{السيد بيضون}. --- نعم، لكن يسبقه \VUgras{النجّار} \textsuperscript{(115)}.

%%K t05-01-25 p
النجّار يصنع \VUgras{هيكل السطح} \textsuperscript{(35)}. يصل \VUgras{العوارض}
\textsuperscript{(37)} بـ\VUgras{الأوتاد} \textsuperscript{(117)}. وأولئك العمال في
الأعلى، بسراويلهم المخملية الواسعة، نجّارون. أحدهم يمسك عارضة صغيرة
والآخر ينحت وتدًا بـ\VUgras{فأسه} \textsuperscript{(116)}. والذي قرب
\VUgras{المِدخنة} \textsuperscript{(41)} هو المُقرمِد. يسمّر الألواح على
\VUgras{الروافد} (*)، و\VUgras{الأردواز} \textsuperscript{(66)} على الألواح.
وأحيانًا يضع القرميد.

%%K t05-noto-1 noto
\VUnotes{143.2mm}{%
(*) \VUgras{Balk-o} = بالألمانية \textit{Balken}، وبالإنجليزية
\textit{joist}، وبالفرنسية \textit{solive}، وبالإيطالية
\textit{travicello}، وبالروسية \textit{balka}، وبالإسبانية والبرتغالية
\textit{viga}. جذر فني لم يُعتمد بعد، لكنه يُقترح هنا لأداء معنى الكلمة
الفرنسية \textit{solive}، وهي ليست \textit{trabo} (بالفرنسية
\textit{poutre}) ولا عارضة صغيرة. إنها خشبة مربّعة تحمل أرضية وسقفًا.%
}

%%K t05-01-26 p
سطح الإسطبل \VUgras{من القرميد} \textsuperscript{(55)}، وسطح البيت
\VUgras{من الأردواز} \textsuperscript{(32)}، وسطح الشرفة \VUgras{من الزنك}
\textsuperscript{(24)}.

%%K t05-01-27 p
\textsc{يوحنا}. --- وهل يصنع المُقرمِد سطوح الزنك أيضًا؟

%%K t05-01-28 p
\textsc{السيد بيضون}. --- لا، ذاك \VUgras{عامل الزنك} \textsuperscript{(121)} أو
\VUgras{السمكري}، وهو الذي يركّب \VUgras{نافذة السطح} \textsuperscript{(38)} في
سطح البيت. وهو الذي يضع \VUgras{مزاريب السطح} \textsuperscript{(43)}
و\VUgras{أنابيب النزول} \textsuperscript{(42)}. يلحم الزنك والصفيح. وهو الذي
ثبّت \VUgras{مانعة الصواعق} \textsuperscript{(40)} و\VUgras{دوّارة الريح}
\textsuperscript{(39)} على \VUgras{الجملون} \textsuperscript{(36)}.

%%K t05-01-29 p
\textsc{يوحنا}. --- فهل تم البيت إذًا؟

%%K t05-tit-5 sub
\VUsaut{3.0mm}
\VUpk{10.2pt}{40}{الأدوات.}
\VUsaut{3.0mm}

%%K t05-01-30 p
\textsc{السيد بيضون}. --- ليس تمامًا. \VUgras{المُلبِّس} \textsuperscript{(99)}
يبني بـ\VUgras{الطوب} \textsuperscript{(61)} الحواجز التي تقسمه إلى غرف مختلفة.
ويكسو \VUgras{الأسقف} \textsuperscript{(26)} والجدران بـ\VUgras{الجص}
\textsuperscript{(70)}. وهو الآن يعمل في سقف \VUgras{الشرفة} \textsuperscript{(33)}.
وله، كالبنّاء، \VUgras{مالج} \textsuperscript{(100)} و\VUgras{حوض}
\textsuperscript{(101)}.

%%K t05-01-31 p
ثم يصنع النجّار الأبواب والنوافذ و\VUgras{الأرضيات الخشبية}
\textsuperscript{(16)} و\VUgras{المصاريع} \textsuperscript{(19)}.

%%K t05-01-32 p
وهو الآن يعمل في الفناء على \VUgras{منضدته} \textsuperscript{(90)}. له
\VUgras{منشار} \textsuperscript{(94)} لنشر \VUgras{الخشب} \textsuperscript{(69)}،
و\VUgras{مسحاج} \textsuperscript{(91)}، و\VUgras{ملزمة} \textsuperscript{(92)}،
و\VUgras{إزميل} \textsuperscript{(94 \textit{bis})}، و\VUgras{فِرجار}
\textsuperscript{(95)}، و\VUgras{مطرقة خشبية} \textsuperscript{(95 \textit{bis})}.

%%K t05-01-33 p
\textsc{يوحنا}. --- آه، إن كون المرء نجّارًا أمتع من كونه بنّاءً! يا عمّي،
أتشتري لي منضدة ومنشارًا ومسحاجًا؟ سأصنع بيتًا لميدور.

%%K t05-01-34 p
\textsc{العم}. --- نعم يا بنيّ.

%%K t05-01-35 p
\textsc{يوحنا}. --- ما أسعدني! ومن ذلك الرجل الواقف قرب الباب الأمامي؟

%%K t05-tit-6 sub
\VUsaut{3.0mm}
\VUpk{10.2pt}{40}{الدهان.}
\VUsaut{3.0mm}

%%K t05-01-36 p
\textsc{السيد بيضون}. --- إنه \VUgras{دهّان} \textsuperscript{(102)}. يقف على
سلّمه ومعه علبة \VUgras{دهان} \textsuperscript{(103)} و\VUgras{فرشاته}
\textsuperscript{(104)}. يطلي خشب الباب الأمامي ليدوم أطول.

%%K t05-01-37 p
وهو أيضًا الذي يلصق ورق الجدران في الغرف.

%%K t05-01-38 p
و\VUgras{المنجّد} \textsuperscript{(107)} على \VUgras{سلّم} \textsuperscript{(106)}،
عند \VUgras{الشرفة} \textsuperscript{(28)}، يركّب \VUgras{ستارة} \textsuperscript{(29)}.

%%K t05-01-39 p
والعامل الواقف قرب النافذة الكبيرة هو \VUgras{الزجّاج} \textsuperscript{(96)}.
يثبّت \VUgras{الألواح} \textsuperscript{(18)} بـ\VUgras{المعجون} \textsuperscript{(73)}.
وذلك العامل الآخر، الجاثي قرب باب القبو، \VUgras{حدّاد أقفال}
\textsuperscript{(97)}. يركّب \VUgras{قفلًا} \textsuperscript{(6)}. و\VUgras{مِبرده}
\textsuperscript{(98)} إلى جانبه.

%%K t05-01-40 p
\textsc{يوحنا}. --- وهل الرجل الذي يضرب قطعة حديد بـ\VUgras{مطرقته}
\textsuperscript{(84)} حدّاد أقفال أيضًا؟

%%K t05-01-41 p
\textsc{السيد بيضون}. --- بل هو، على الأدق، \VUgras{حدّاد} \textsuperscript{(125)}.
يطرق \VUgras{حديدًا} \textsuperscript{(67)} لـ\VUgras{الكهربائي}
\textsuperscript{(124)} الذي هو فوق سطح \VUgras{المرآب} \textsuperscript{(51)}. وله
\VUgras{كور} \textsuperscript{(126)} و\VUgras{سندان} \textsuperscript{(127)}
و\VUgras{كلّابة} \textsuperscript{(128)}.

%%K t05-01-42 p
والكهربائي يثبّت \VUgras{السلك النحاسي} \textsuperscript{(44)} للهاتف.

%%K t05-tit-7 sub
\VUpk{10.2pt}{40}{أعمال البستنة.}
\VUsaut{3.0mm}

%%K t05-01-43 p
\textsc{العم}. --- في آخر \VUgras{الفناء} \textsuperscript{(45)}، قرب
\VUgras{السياج الحديدي} \textsuperscript{(47)} و\VUgras{بوابة العربات}
\textsuperscript{(48)}، ترى \VUgras{البستاني} \textsuperscript{(129)}. إنه يهيّئ
\VUgras{الحديقة} \textsuperscript{(49)} الصغيرة. قلب التربة بـ\VUgras{مِسحاته}
\textsuperscript{(131)}، وهو يبذر ويسوّي بـ\VUgras{المِجرفة المشطية}
\textsuperscript{(130)}. ثم يسقي بـ\VUgras{مِرشّته} \textsuperscript{(132)}
\VUgras{الأحواض} \textsuperscript{(150)} فتنمو النباتات.

%%K t05-01-44 p
و\VUgras{السائس} \textsuperscript{(133)}، الذي فرغ لتوّه من تنظيف الحصان وعلفه،
ينظف \VUgras{العربة} \textsuperscript{(52)} بإسفنجة و\VUgras{مضخة يدوية}
\textsuperscript{(134)} ودلو ماء. ثم يدخلها المرآب ويغلق الباب
بـ\VUgras{المزلاج} \textsuperscript{(53)} الثقيل.

%%K t05-01-45 p
ها قد اكتفيت الآن، على ما أرجو؛ فقد نلت من الشرح ما يكفي.

%%K t05-01-46 p
\textsc{يوحنا}. --- نعم يا عمّي، لكني أودّ أن أرى داخل البيت.

%%K t05-01-47 p
\textsc{العم}. --- ذلك يكون غدًا. سيشرح لك السيد رو المخطط ويسمّي لك كل
غرفة.

%%K t05-tit-8 sub
\VUsaut{3.0mm}
\VUpk{10.2pt}{40}{تقسيم البيت.}
\VUsaut{2.4mm}

%%K t05-apar-2 apar
{\centering\textit{(انظر المخطط.)}\par}
\VUsaut{2.4mm}

%%K t05-01-48 p
\textsc{المهندس}. --- تعال يا يوحنا، سأطوف بك في أجزاء البيت المختلفة.

%%K t05-01-49 p
لندخل \VUgras{الطابق الأرضي} \textsuperscript{(7)}. هذا زرّ \VUgras{الجرس}
\textsuperscript{(11)}. نصعد \VUgras{الدرجات} \textsuperscript{(14)} الثلاث
لـ\VUgras{السلّم الأمامي} \textsuperscript{(9)}، ونجتاز \VUgras{عتبة}
\textsuperscript{(10)} \VUgras{الباب الأمامي} \textsuperscript{(8)}، فإذا نحن عند
أسفل \VUgras{الدرج} \textsuperscript{(13)}.

%%K t05-01-50 p
\textsc{يوحنا}. --- هذا هو الممر.

%%K t05-01-51 p
\textsc{المهندس}. --- لا، هذه \VUgras{الردهة} \textsuperscript{(12)}.
و\VUgras{الممر} \textsuperscript{(a)} في الآخر. وعلى اليمين \VUgras{صالة
الاستقبال} \textsuperscript{(b)} و\VUgras{غرفة الطعام} \textsuperscript{(c)}؛ وقُبالة
غرفة الطعام \VUgras{المطبخ} \textsuperscript{(d)} و\VUgras{المُؤونة}
\textsuperscript{(e)}؛ ثم، من جهة الواجهة، \VUgras{المكتب} \textsuperscript{(f)}.
و\VUgras{الشرفة} \textsuperscript{(23)}، بـ\VUgras{بابها الزجاجي}
\textsuperscript{(25)}، تجاور صالة الاستقبال.

%%K t05-01-52 p
والآن نصعد الدرج. أمسك \VUgras{الدرابزين} \textsuperscript{(15)} وعُدّ
الدرجات. إنها أربع عشرة. وهذا \VUgras{البسطة} \textsuperscript{(m)}: نحن في
\VUgras{الطابق} \textsuperscript{(27)} الأول.

%%K t05-tit-9 sub
\VUsaut{3.0mm}
\VUpk{10.2pt}{40}{الطابق الأول.}
\VUsaut{3.0mm}

%%K t05-01-53 p
قُبالة الدرج \VUgras{المرحاض} \textsuperscript{(20)} و\VUgras{غرفة الزينة}
\textsuperscript{(j)}. وعلى اليمين \VUgras{غرفة نوم} \textsuperscript{(g)} جميلة لها
نافذتان على \VUgras{واجهة} \textsuperscript{(1)} البيت. وعلى اليسار
\VUgras{الحمّام} \textsuperscript{(1)} و\VUgras{غرفة الضيوف} \textsuperscript{(i)} وفيها
\VUgras{مِخدع} \textsuperscript{(h)} كبير. وإلى جانب غرفة النوم \VUgras{غرفة
الأطفال} \textsuperscript{(k)}. وهي تُفضي إلى \VUgras{سطح} \textsuperscript{(21)}
واسع. وتحت هذا السطح مرحاض آخر \textsuperscript{(20)}، وعلى الجدار ها هو
\VUgras{الجرس} \textsuperscript{(22)} الذي يُقرَع للعشاء.

%%K t05-01-54 p
\textsc{يوحنا}. --- لنصعد إلى \VUgras{الطابق الثاني}.

%%K t05-01-55 p
\textsc{المهندس}. --- ليس ثمة طابق ثانٍ. تحت السطح \VUgras{غرف علوية}
\textsuperscript{(33)} والعِلّية \textsuperscript{(34)} فحسب. فرغنا من جولتنا،
ويمكننا النزول.

%%K t05-01-56 p
\textsc{يوحنا}. --- شكرًا لك يا سيدي، كان هذا ممتعًا جدًا. سأحكي لأبي كل
ما رأيت.
\VUsaut{4.0mm}
\VUfilet{22mm}
